\documentclass[12pt]{article}
\usepackage{geometry}
\usepackage{amsmath,amssymb}
\usepackage{graphicx}
\geometry{margin=1in}

\begin{document}

\title{Module SAT.O$_{10}$ (\texttt{sat\_o10.py})}
\author{SAT.4D Coherence Project}
\date{\today}
\maketitle

\section*{Overview}
Provides the two-sector calibration and prediction framework:
\begin{itemize}
  \item \emph{Leptonic sector}: fit \(m^2 = p_2\,Q^2 + p_1\) to \((e,\mu,\pi)\)
  \item \emph{Hadronic sector}: fit the same form to \((\rho,K^*,\phi)\)
  \item Predict ground states and radial excitations
  \item Plot trajectories with data overlays
\end{itemize}

\section{Functions}

\subsection{\texttt{calibrate\_leptonic}}
\begin{verbatim}
def calibrate_leptonic(masses, Q_vals):
    # Fit m^2 = p2*Q^2 + p1 to leptonic ground states
    # Returns {'p2':..., 'p1':..., 'masses':..., 'Q_vals':...}
    m2 = np.array(masses)**2
    A = np.vstack([np.array(Q_vals)**2, np.ones(len(Q_vals))]).T
    p2, p1 = np.linalg.lstsq(A, m2, rcond=None)[0]
    return {'p2': p2, 'p1': p1, 'masses': masses, 'Q_vals': Q_vals}
\end{verbatim}

\subsection{\texttt{calibrate\_hadronic}}
\begin{verbatim}
def calibrate_hadronic(masses, Q_vals):
    # Fit m^2 = p2*Q^2 + p1 to hadronic ground states
    # Returns {'p2':..., 'p1':..., 'masses':..., 'Q_vals':...}
    m2 = np.array(masses)**2
    A = np.vstack([np.array(Q_vals)**2, np.ones(len(Q_vals))]).T
    p2, p1 = np.linalg.lstsq(A, m2, rcond=None)[0]
    return {'p2': p2, 'p1': p1, 'masses': masses, 'Q_vals': Q_vals}
\end{verbatim}

\subsection{\texttt{predict\_ground\_states}}
\begin{verbatim}
def predict_ground_states(p2, p1, Q_range):
    # Compute m_ground(Q)=sqrt(p2*Q^2 + p1) for Q in Q_range
    records = []
    for Q in Q_range:
        m2 = p2 * Q**2 + p1
        m = np.sqrt(m2) if m2 > 0 else float('nan')
        records.append({'Q': Q, 'm_ground': m})
    return pd.DataFrame(records)
\end{verbatim}

\subsection{\texttt{predict\_radial\_states}}
\begin{verbatim}
def predict_radial_states(p2, p1, k_rad, Q_vals):
    # Compute m_radial(Q)=sqrt(p2*Q^2 + p1 + k_rad) for Q in Q_vals
    records = []
    for Q in Q_vals:
        m2 = p2 * Q**2 + p1 + k_rad
        m = np.sqrt(m2) if m2 > 0 else float('nan')
        records.append({'Q': Q, 'm_radial': m})
    return pd.DataFrame(records)
\end{verbatim}

\subsection{\texttt{plot\_trajectories}}
\begin{verbatim}
def plot_trajectories(params_lept, params_had):
    # Plot m^2 vs Q² for both sectors with data points
    Ql = np.linspace(min(params_lept['Q_vals']), max(params_lept['Q_vals'])*1.5, 200)
    m2l = params_lept['p2'] * Ql**2 + params_lept['p1']
    Qh = np.linspace(min(params_had['Q_vals']), max(params_had['Q_vals'])*1.5, 200)
    m2h = params_had['p2'] * Qh**2 + params_had['p1']
    plt.figure()
    plt.plot(Ql, m2l, label='Leptonic')
    plt.scatter(params_lept['Q_vals'], np.array(params_lept['masses'])**2, marker='o')
    plt.plot(Qh, m2h, label='Hadronic')
    plt.scatter(params_had['Q_vals'], np.array(params_had['masses'])**2, marker='s')
    plt.xlabel('Q')
    plt.ylabel(r'$m^2$ (GeV$^2$)')
    plt.legend()
    plt.show()
\end{verbatim}

\section*{Usage Example}
\begin{verbatim}
from sat_o.sat_o10 import (
    calibrate_leptonic, calibrate_hadronic,
    predict_ground_states, predict_radial_states,
    plot_trajectories
)

lep_params = calibrate_leptonic(
    masses=[0.000511,0.10566,0.13957], Q_vals=[1,2,3]
)
had_params = calibrate_hadronic(
    masses=[0.775,0.892,1.019], Q_vals=[3,4,6]
)
df_lep = predict_ground_states(**lep_params, Q_range=range(1,11))
k_rad = 1.450**2 - (had_params['p2']*3**2 + had_params['p1'])
df_rad = predict_radial_states(**had_params, k_rad=k_rad, Q_vals=[3,4,6])
plot_trajectories(lep_params, had_params)
\end{verbatim}

\end{document}
